\documentclass[../../main.tex]{subfiles}
\begin{document}

\begin{center}
  {\fontsize{48pt}{58pt}\selectfont \textbf{2004}\normalsize\textbf{年度}}
\end{center}

\begin{table}[htb]
  \centering
  \caption{出題テーマと難易度}
  \begin{tabular}{|c|p{9cm}|c|}
    \hline
    \textbf{問題番号} & \textbf{テーマ} & \textbf{難易度} \\
    \hline
    第1問 & 三角関数の最大値・最小値 & \\
    \hline
    第2問 & 対数関数のグラフと面積 & \\
    \hline
    第3問 & 多項式の余りの極限 & \\
    \hline
    第4問 & 行列方程式の全射条件 & \\
    \hline
    第5問 & 円を通る共役複素数の関係 & \\
    \hline
    第6問 & 箱の玉を交換する確率 & \\
    \hline
  \end{tabular}
\end{table}

\begin{table}[htb]
  \centering
  \caption{解答}
  \renewcommand{\arraystretch}{1.4}
  \begin{tabular}{|c|c|p{9cm}|}
    \hline
    \textbf{問題番号} & \textbf{小問} & \textbf{答え} \\
    \hline
    第1問 & -- & 最大値$1$，最小値$-7/2$ \\
    \hline
    第2問 & -- & $\dfrac{e-5/2}{a^2}$ \\
    \hline
    第3問 & -- & $\lim_{n\to\infty}a_n=e^{-2}$，$\lim_{n\to\infty}b_n=0$ \\
    \hline
    第4問 & -- &  \\
    \hline
    第5問 & -- & $C$は$-1/\bar{\alpha}$も通る（証明問題） \\
    \hline
    第6問 & -- & $\dfrac{1}{N}\left\{1-\left(\dfrac{N-1}{N}\right)^N\right\}$ \\
    \hline
  \end{tabular}
\end{table}

\paragraph{試験形態}
6問．試験時間は要確認．

\end{document}
